\documentclass[10pt,a4paper]{amsart}
\usepackage[margin=19mm]{geometry}
\usepackage{amsmath,amssymb,mathtools}
\usepackage[scr=rsfs,cal=euler]{mathalfa}
\usepackage{xfrac}
\usepackage[unicode,hidelinks,bookmarks=false]{hyperref}
\newcommand{\ie}{i.e.\ }
\newcommand{\cf}{cf.\ }
\newcommand{\resp}{respectively\ }
\newcommand{\Subsection}{\subsection}
\providecommand{\nobreakdash}{\nobreak\hbox{-}}
\setlength{\emergencystretch}{2em}
\multlinegap0pt
\usepackage{amssymb}
\usepackage{mathtools}
\usepackage{xcolor}
\usepackage[shortlabels]{enumitem}
\usepackage{tikz}
\newtheorem{lemma}{Lemma}
\newcommand{\C}{\mathbb{C}}
\newcommand{\Ev}{\mathcal{E}}
\newcommand{\Mm}{\mathcal{M}}
\newcommand{\Om}{\Omega}
\newcommand{\R}{\mathbb{R}}
\newcommand{\abs}[1]{\left\lvert#1\right\rvert}
\title{Endogenous Feedback in Size-Structured Transport Equations}
\author{Jiguang Yu, Louis Shuo Wang}
\date{Source: arXiv:2607.02877v2 (CC BY 4.0)}
\begin{document}
\maketitle

\noindent\emph{Source and licence.} Endogenous Feedback in Size-Structured Transport Equations. Version arXiv:2607.02877v2 (CC BY 4.0), distributed by arXiv under the Creative Commons Attribution 4.0 International licence, http://creativecommons.org/licenses/by/4.0/, as recorded in the arXiv licence metadata for this exact version and shown on the abstract page https://arxiv.org/abs/2607.02877v2. Full licence notice: \texttt{licenses/arxiv-2607.02877-CC-BY-4.0.txt}.\par\smallskip

\noindent\textit{Editorial scope.} Counted unit: the article's lemma lem:main-Ev-props (source0.tex lines 1050--1085). The article's complete original proof of it is delivered with it (source0.tex lines 2309--2563, one \texttt{proof} environment of 255 lines, which the article places away from the statement). No mathematics is written, repaired or re-derived: the statement and the proof are the article's own lines, in the article's own notation, and no retained line is substituted or re-flowed.  No further source-local context is needed: every cross-reference of every retained line resolves inside the two delivered spans.  Nothing was omitted from any retained span, so the package carries no omission record.  No retained line contains a citation, so the wrapper carries no bibliography and no bibliography entry.  No retained line contains a figure environment, an image-inclusion command or drawing code, so no image is delivered and the package declares no asset dependency.  Editorial formatting only, in addition: an AMS \texttt{amsart} preamble that restates the article's own declarations and macros used by the retained lines, namely the environments \texttt{lemma} on separate counters, together with the standard packages those lines need.  The retained environments print in this document as Lemma 1 in order of appearance, whereas the article prints the same items with the numbering of its own full document.  The complete native source of the article as distributed in the arXiv e-print for this exact version is bundled unchanged as \texttt{sources/204346/source0.tex}.
\begin{lemma}[Properties of \(\Ev\)]
\label{lem:main-Ev-props}
The map \(\Ev:\C\to\C\) is entire,
\[
\Ev(\bar\lambda)=\overline{\Ev(\lambda)},
\qquad
\Ev(\lambda)=\int_0^{t_\sharp}e^{-\lambda\tau}\kappa(\tau)\,d\tau .
\]
Moreover,
$\Ev(\lambda)\to0$ as $\lambda\to+\infty$, $\lambda\in\R$,
and, for \(\Re\lambda\ge0\),
$\displaystyle 
\abs{\Ev(\lambda)}\le G_0$,
where
\[
G_0:=
\int_{l_0}^{l_m}
\frac{\chi(l)}{g^*(l)}
\Pi^0(l)
\left[
\abs{a_0}
+
\int_{l_0}^{l}
\Pi^0(s)^{-1}\,d\abs{\sigma_{\rm int}}(s)
\right]dl,
\]
with
$\displaystyle 
\Pi^0(l):=
\exp\!\left[
-\int_{l_0}^{l}\frac{\mu^*(\xi)+u(\xi)}{g^*(\xi)}\,d\xi
\right]$.
Finally,
$\displaystyle 
\Ev(0)=\Phi'(E^*)$.
\end{lemma}

\begin{proof}[Proof of Lemma~\ref{lem:main-Ev-props}]
Let \(\lambda\in\C\). The inhomogeneous stationary transport problem is
\[
\frac{d}{dl}\big(g^*\varphi_1^\lambda\big)
+
(\lambda+\mu^*+u)\varphi_1^\lambda
=
\sigma_{\rm int},
\qquad
g^*(l_0)\varphi_1^\lambda(l_0)=-a_0 .
\]
Set
\[
Y^\lambda(l):=g^*(l)\varphi_1^\lambda(l),
\qquad
Q^\lambda(l):=
\frac{\lambda+\mu^*(l)+u(l)}{g^*(l)}.
\]
Then
$\displaystyle 
(Y^\lambda)'(l)+Q^\lambda(l)Y^\lambda(l)=\sigma_{\rm int}$
in the distributional sense, and
$\displaystyle 
Y^\lambda(l_0)=-a_0$.
Define
\[
\Pi^\lambda(l)
:=
\exp\!\left[-\int_{l_0}^{l}Q^\lambda(\xi)\,d\xi\right]
=
\exp\!\left[-\int_{l_0}^{l}
\frac{\lambda+\mu^*(\xi)+u(\xi)}{g^*(\xi)}\,d\xi
\right].
\]
Then
$\displaystyle 
\frac{d}{dl}\left((\Pi^\lambda(l))^{-1}Y^\lambda(l)\right)
=
(\Pi^\lambda(l))^{-1}\,d\sigma_{\rm int}(l)$.
Therefore
\[
(\Pi^\lambda(l))^{-1}Y^\lambda(l)
=
Y^\lambda(l_0)
+
\int_{l_0}^{l}(\Pi^\lambda(s))^{-1}\,d\sigma_{\rm int}(s),
\]
and hence
\[
Y^\lambda(l)
=
\Pi^\lambda(l)
\left[
-a_0+\int_{l_0}^{l}(\Pi^\lambda(s))^{-1}\,d\sigma_{\rm int}(s)
\right].
\]
Since \(Y^\lambda=g^*\varphi_1^\lambda\),
\[
\varphi_1^\lambda(l)
=
\frac{\Pi^\lambda(l)}{g^*(l)}
\left[
-a_0+\int_{l_0}^{l}(\Pi^\lambda(s))^{-1}\,d\sigma_{\rm int}(s)
\right].
\]
Thus
\begin{equation}
\label{eq:proof-Ev-explicit}
\Ev(\lambda)
=
\int_{l_0}^{l_m}
\frac{\chi(l)}{g^*(l)}
\Pi^\lambda(l)
\left[
-a_0+\int_{l_0}^{l}(\Pi^\lambda(s))^{-1}\,d\sigma_{\rm int}(s)
\right]dl .
\end{equation}
For each \(l\),
$\displaystyle 
\Pi^\lambda(l)
=
\exp\!\left[-\lambda\int_{l_0}^{l}\frac{d\xi}{g^*(\xi)}\right]
\Pi^0(l)$,
hence \(\lambda\mapsto\Pi^\lambda(l)\) is entire. Likewise,
$\displaystyle 
(\Pi^\lambda(s))^{-1}
=
\exp\!\left[\lambda\int_{l_0}^{s}\frac{d\xi}{g^*(\xi)}\right]
(\Pi^0(s))^{-1}$
is entire. For \(\lambda\) in a compact set \(K\subset\C\),
\[
\abs{\Pi^\lambda(l)}
\le
e^{C_Kt_\sharp}\Pi^0(l),
\qquad
\abs{(\Pi^\lambda(s))^{-1}}
\le
e^{C_Kt_\sharp}(\Pi^0(s))^{-1},
\]
where
$\displaystyle 
C_K:=\sup_{\lambda\in K}\abs{\lambda}$.
Since \(\sigma_{\rm int}\in\Mm(\Om)\), the integrand in
\eqref{eq:proof-Ev-explicit} is locally dominated uniformly in \(\lambda\). Hence
$\Ev: \mathbb C\mapsto \mathbb C$ is entire.
Because all coefficients and measures are real,
\[
\overline{\Pi^\lambda(l)}=\Pi^{\bar\lambda}(l),
\qquad
\overline{\int_{l_0}^{l}(\Pi^\lambda(s))^{-1}\,d\sigma_{\rm int}(s)}
=
\int_{l_0}^{l}(\Pi^{\bar\lambda}(s))^{-1}\,d\sigma_{\rm int}(s),
\]
so
$\displaystyle 
\Ev(\bar\lambda)=\overline{\Ev(\lambda)}$.
Next, compare \(\Ev\) with the Laplace transform of the renewal kernel. Let
$v(t,l)=e^{\lambda t}\varphi(l)$ and 
$e(t)=e^{\lambda t}\hat e$.
Substitution into the linearized system gives
\[
\frac{d}{dl}(g^*\varphi)+(\lambda+\mu^*+u)\varphi
=
\sigma_{\rm int}\hat e,
\qquad
g^*(l_0)\varphi(l_0)=-a_0\hat e .
\]
Thus
$\displaystyle 
\varphi=\hat e\,\varphi_1^\lambda$.
The consistency condition \(e=\langle\chi,v\rangle\) gives
$\displaystyle 
\hat e
=
\langle\chi,\varphi\rangle
=
\hat e\,\langle\chi,\varphi_1^\lambda\rangle
=
\hat e\,\Ev(\lambda)$.
On the other hand, for a homogeneous modal solution of the scalar renewal
equation,
\[
e(t)=e^{\lambda t}\hat e,
\qquad
f=0,
\]
we have
$\displaystyle 
e^{\lambda t}\hat e
=
\int_0^t\kappa(t-s)e^{\lambda s}\hat e\,ds$.
Writing \(\tau=t-s\),
$\displaystyle 
e^{\lambda t}\hat e
=
e^{\lambda t}\hat e\int_0^t e^{-\lambda\tau}\kappa(\tau)\,d\tau$.
Since \(\operatorname{supp}\kappa\subset[0,t_\sharp]\), for \(t\ge t_\sharp\),
$\displaystyle 
\hat e
=
\hat e\int_0^{t_\sharp}e^{-\lambda\tau}\kappa(\tau)\,d\tau$.
Therefore the two modal consistency factors coincide:
\[
\Ev(\lambda)
=
\int_0^{t_\sharp}e^{-\lambda\tau}\kappa(\tau)\,d\tau.
\]
Now let \(\lambda\to+\infty\), \(\lambda\in\R\). In \eqref{eq:proof-Ev-explicit},
\[
\Pi^\lambda(l)
=
e^{-\lambda\tau(l)}\Pi^0(l),
\qquad
\frac{\Pi^\lambda(l)}{\Pi^\lambda(s)}
=
e^{-\lambda(\tau(l)-\tau(s))}
\frac{\Pi^0(l)}{\Pi^0(s)}.
\]
For \(l>l_0\),
$\displaystyle 
e^{-\lambda\tau(l)}\to0$.
For \(s<l\),
$\displaystyle  e^{-\lambda(\tau(l)-\tau(s))}\to0$.
The diagonal \(s=l\) is irrelevant for the absolutely continuous part and is
controlled for atoms by the one-sided Stieltjes convention in the variation-of-constants formula. The dominated bounds
\[
0\le e^{-\lambda\tau(l)}\le1,
\qquad
0\le e^{-\lambda(\tau(l)-\tau(s))}\le1
\]
give
$\displaystyle 
\Ev(\lambda)\to0$.
For \(\Re\lambda\ge0\),
$\displaystyle 
\abs{\Pi^\lambda(l)}
=
e^{-\Re\lambda\,\tau(l)}\Pi^0(l)
\le
\Pi^0(l)$,
and for \(s\le l\),
$\displaystyle 
\frac{\abs{\Pi^\lambda(l)}}{\abs{\Pi^\lambda(s)}}
=
e^{-\Re\lambda(\tau(l)-\tau(s))}
\frac{\Pi^0(l)}{\Pi^0(s)}
\le
\frac{\Pi^0(l)}{\Pi^0(s)}$.
Therefore, from \eqref{eq:proof-Ev-explicit},
\[
\abs{\Ev(\lambda)}
\le
\int_{l_0}^{l_m}
\frac{\chi(l)}{g^*(l)}
\Pi^0(l)
\left[
\abs{a_0}
+
\int_{l_0}^{l}(\Pi^0(s))^{-1}\,d\abs{\sigma_{\rm int}}(s)
\right]dl
=
G_0.
\]
Finally, put \(\lambda=0\). The equation for \(\varphi_1^0\) is
\[
\frac{d}{dl}(g^*\varphi_1^0)+(\mu^*+u)\varphi_1^0
=
\sigma_{\rm int},
\qquad
g^*(l_0)\varphi_1^0(l_0)=-a_0.
\]
The forward sensitivity
$\displaystyle 
y(l):=\partial_E x_E(l)|_{E=E^*}$
satisfies exactly
\[
\frac{d}{dl}(g^*y)+(\mu^*+u)y
=
\sigma_{\rm int},
\qquad
g^*(l_0)y(l_0)=-a_0 .
\]
Uniqueness of the first-order boundary-value problem gives
$\displaystyle 
\varphi_1^0=y$.
Therefore
$\displaystyle 
\Ev(0)
=
\langle\chi,\varphi_1^0\rangle
=
\langle\chi,y\rangle
=
\Phi'(E^*)$.
\end{proof}

\end{document}
